% Deep learning related work (Harshil)
% Complements Section sec:related (Priyanshu), which covers the NDWS spread models. This section covers the
% methods behind the Phase 2 design choices. Our numbers come from reports/dl/results_dl.csv,
% reports/dl/checks/aug_ablation.md, reports/dl/hybrid_ensemble_summary.md and
% reports/dl/architecture_ablation_summary.md.

\section{Deep Learning Related Work}
\label{sec:dl_related}

Section~\ref{sec:related} reviewed the spread models trained on NDWS. Here we review the methods behind
the choices of our Phase~2 model: the architecture and how much of the patch it sees, the loss, the
augmentation, the hybrid with the Phase~1 model, and training. Table~\ref{tab:dl_relwork} lists what each
paper contributes and what we found when we tested it.

\textbf{Long et al.~\cite{long2015}} turn classification networks into fully convolutional networks (FCNs)
for semantic segmentation. Replacing the fully connected layers by convolutions lets one network take an
input of any size and output a dense map of the same size, and skip connections combine coarse, deep
features with fine, shallow ones. \relevance{} all our models are fully convolutional, which is what lets
a model trained on $32\times32$ crops be scored on $64\times64$ patches (Section~\ref{sec:dl_theory}).

\textbf{Ronneberger et al.~\cite{ronneberger2015}} (U-Net) make the FCN symmetric: an encoder halves the
resolution at each level, a decoder upsamples back, and each decoder level concatenates the encoder
features of the same resolution. \relevance{} U-Net is our main model. Its depth, the number of
levels, sets how far each output pixel can see.

\textbf{Luo et al.~\cite{luo2016}} show that the \emph{effective} receptive field of a CNN, the region
that actually moves an output, is roughly Gaussian and covers only a fraction of the theoretical receptive
field. It grows only with the square root of the number of stacked convolutions, while downsampling
enlarges it much faster. \relevance{} this predicts our architecture ablation: an extra U-Net level (one
more downsampling) raised test PR AUC by 0.064, while doubling the width at a similar parameter count did
not help.

\textbf{Shi et al.~\cite{shi2015}} (ConvLSTM) replace the matrix products of an LSTM by convolutions, so a
recurrent model keeps spatial structure, and use it to forecast rainfall from sequences of radar maps.
\relevance{} ConvLSTM needs several input days. NDWS gives one day of inputs per sample, so we do not use
it (multi-day data such as WildfireSpreadTS~\cite{gerard2023} would allow it).

\textbf{Cohen and Welling~\cite{cohen2016}} build group-equivariant CNNs whose layers commute with
$90^\circ$ rotations and reflections, so rotating the input rotates the output exactly. Data augmentation
with the same transforms only \emph{encourages} this behaviour. \relevance{} we tested the augmentation
route (Section~\ref{sec:dl_aug}). On our data it did not help, and a model that is equivariant by
construction would bake in a symmetry the data do not have (Section~\ref{sec:dl_theory}).

\textbf{Wolpert~\cite{wolpert1992}} introduces stacked generalisation: a second-level model learns from
the predictions of first-level models, and those predictions must come from models that did not see the
sample, usually via cross-validation, or the second level learns to trust overfitted inputs.
\relevance{} our hybrid U-Net takes the XGBoost probability as a 13th channel, computed out of fold with
folds grouped by fire event. It changed test PR AUC by less than 0.003.

\textbf{Prechelt~\cite{prechelt1998}} compares rules for stopping training when the validation error
stops falling. Waiting longer before stopping gains a little generalisation at a large cost in training
time. \relevance{} we stop when eval PR AUC has not improved for 5 epochs and keep the best epoch.

\textbf{Focal loss}~\cite{lin2017} (Section~\ref{sec:dl_loss}) and batch normalisation and
Adam~\cite{ioffe2015,kingma2015} (Section~\ref{sec:dl_theory}) complete the training setup.

\begin{table}[t]
\centering
\caption{Methods behind the Phase~2 model and what we found when we tested them. PR AUC is on the test
split, mean $\pm$ standard deviation over three seeds unless marked as a single run.}
\label{tab:dl_relwork}
\footnotesize
\setlength{\tabcolsep}{3pt}
\renewcommand{\arraystretch}{1.12}
\newcolumntype{L}[1]{>{\raggedright\arraybackslash}p{#1}}
\begin{tabular}{L{1.95cm}L{2.45cm}L{3.8cm}}
\toprule
\rowcolor{tblhead}
\textbf{Paper} & \textbf{Idea} & \textbf{Our test and result} \\
\midrule
FCN~\cite{long2015} & dense prediction, any input size & train on centre $32\times32$, score on $64\times64$: 0.172 against 0.269 (single run) \\
U-Net~\cite{ronneberger2015} & encoder-decoder with skips & backbone; CNN without downsampling ties XGBoost \\
ERF~\cite{luo2016} & downsampling grows the effective field & depth 3 $\to$ 4: 0.280 $\pm$ 0.002 $\to$ 0.344 $\pm$ 0.007; width 64: no gain \\
Focal loss~\cite{lin2017} & down-weight easy pixels & about $+0.007$ over weighted BCE; kept \\
G-CNN~\cite{cohen2016} & rotation equivariance & flips/rotations with wind turned: no gain; not used \\
Stacking~\cite{wolpert1992} & out-of-fold inputs & XGBoost channel: $+0.002$ / $-0.003$ (single runs); not used \\
ConvLSTM~\cite{shi2015} & recurrent spatial model & not applicable: one input day \\
Early stopping~\cite{prechelt1998} & stop on validation & eval PR AUC, patience 5; best epochs 9--16 \\
\bottomrule
\end{tabular}
\end{table}

\begin{projectbox}
\textbf{Where Phase~2 sits.} Our final U-Net (depth 4, focal loss) reaches 0.344 $\pm$ 0.007 on the
official split with full patches. That is above Huot et al.'s CNN (0.284) and in line with the U-Net of
Esparza et al.~\cite{esparza2026} (0.346), the only other result we found under the same protocol. The
methods that changed the score were the ones that changed how much of the patch the model sees. Loss,
augmentation and the Phase~1 features moved it by less than the noise between seeds.
\end{projectbox}
